\documentclass[10pt]{article}
\usepackage[T1]{fontenc}
\usepackage[utf8]{inputenc}
\usepackage[french]{babel}
\usepackage{amsmath,amssymb,mathrsfs,mathpazo}
\usepackage{geometry}
\geometry{a4paper,margin=2cm}
\usepackage{graphicx,array,tabularx,tabularray,colortbl,multirow,multicol,booktabs,enumitem,pifont,ulem,fancybox,diagbox,cancel}
\usepackage{pstricks,pstricks-add,pst-plot,pst-tree,pst-node,pst-text,pst-func,pst-math,pst-bspline,pst-3dplot,pst-eucl,pst-solides3d}
\let\pscalculate\undefined
\usepackage{pst-bezier}
\usepackage{tikz,pgfplots,tkz-tab}
\usetikzlibrary{arrows,arrows.meta,patterns,calc,positioning,trees,decorations.markings,shapes.geometric,shapes.misc,shapes.arrows,intersections,decorations.pathmorphing,decorations.pathreplacing,backgrounds,shadows,babel}
\pgfplotsset{compat=1.18}
\usepgfplotslibrary{fillbetween}
\usepackage[np]{numprint}
\usepackage{listings,makecell,pas-tableur,twoopt}
\let\barre\undefined
\usepackage[pstricks]{tablvar}
\definecolor{darkgreen}{RGB}{0,100,0}
\providecommand{\red}{\color{red}}
\providecommand{\blue}{\color{blue}}
\providecommand{\green}{\color{green}}
\providecommand{\black}{\color{black}}
\providecommand{\white}{\color{white}}
\providecommand{\gray}{\color{gray}}
\providecommand{\cyan}{\color{cyan}}
\providecommand{\magenta}{\color{magenta}}
\providecommand{\yellow}{\color{yellow}}
\providecommand{\violet}{\color{violet}}
\providecommand{\orange}{\color{orange}}
\providecommand{\euro}{\texteuro}
\providecommand{\EUR}{\texteuro}
\providecommand{\ieme}{\textsuperscript{e}}
\providecommand{\degres}{\ensuremath{^\circ}}
\providecommand{\degree}{\ensuremath{^\circ}}
\providecommand{\vv}[1]{\overrightarrow{#1}}
\providecommand{\Coord}[1]{#1}
\providecommand{\decosix}{}
\providecommand{\starredbullet}{\ensuremath{\bullet}}
\pagestyle{empty}
\setlength{\parindent}{0pt}
\usepackage{adjustbox}\DeclareUnicodeCharacter{2236}{:}\DeclareUnicodeCharacter{2190}{\ensuremath{\leftarrow}}
\providecommand{\R}{\mathbb{R}}
\providecommand{\N}{\mathbb{N}}
\providecommand{\Z}{\mathbb{Z}}
\providecommand{\Q}{\mathbb{Q}}
\providecommand{\D}{\mathbb{D}}
\providecommand{\C}{\mathbb{C}}
\providecommand{\e}{{\mathrm{e}}}
\providecommand{\vect}[1]{\overrightarrow{#1}}
\providecommand{\vectt}[1]{\overrightarrow{\mathrm{#1}}}
\providecommand{\vvt}[1]{\overrightarrow{\mathrm{#1}}}
\providecommand{\np}[2][]{{#2~#1}}
\providecommand{\vv}[1]{{\overrightarrow{#1}}}
\renewcommand{\d}{{\mathrm{d}}}
\providecommand{\degres}{{^\circ}}
\providecommand{\degre}{{^\circ}}
\providecommand{\blue}{\color{blue}}
\providecommand{\red}{\color{red}}
\providecommand{\enskip}{\quad}
\providecommand{\llbracket}{\lbrack\!\lbrack}
\providecommand{\rrbracket}{\rbrack\!\rbrack}
\providecommand{\euro}{{\text{€}}}
\providecommand{\no}{{\text{n°}}}
\providecommand{\interval}[2]{{[#1;#2]}}
\providecommand{\intervFO}[2]{{[#1;#2[}}
\providecommand{\dotfill}{\ldots}

\providecommand{\ds}{\displaystyle}
\providecommand{\textding}[1]{#1}
\providecommand{\Oij}{\(\mathrm{O};\vec{\imath},\vec{\jmath}\)}
\providecommand{\Oijk}{\(\mathrm{O};\vec{\imath},\vec{\jmath},\vec{k}\)}
\providecommand{\Ouv}{\(\mathrm{O};\vec{u},\vec{v}\)}
\providecommand{\intervFF}[2]{\ensuremath{[#1;#2]}}
\renewcommand{\euro}{\ifmmode\text{\texteuro}\else\texteuro\fi}
\DeclareUnicodeCharacter{00B0}{\ifmmode^{\circ}\else\textdegree\fi}
\DeclareUnicodeCharacter{20AC}{\ifmmode\text{\texteuro}\else\texteuro\fi}
% Énoncés : Éducation nationale. Transcription et corrigés : APMEP, Denis Vergès et contributeurs. Découpage et classement : CoopMaths, Sébastien Lozano, Éric Elter et contributeurs. Adaptation de la présentation : Maths Explorer.
% CC BY-SA : https://coopmaths.fr/www/about/
% https://coopmaths.fr/alea/static/bac/2025/tex/bac_2025_09_sujet1_asie_4_cor.tex
\begin{document}
\noindent\textit{APMEP / CoopMaths. CC BY-SA. Présentation adaptée par Maths Explorer.}\par\medskip
Soit $n$ un entier naturel non nul.

Dans le cadre d'une expérience aléatoire, on considère une suite d'évènements $A_{n}$ et on note $p_{n}$ la probabilité de l'évènement $A_{n}$.

\begin{list}{\textbullet}{Pour les parties \textbf{A} et \textbf{B} de l'exercice, on considère que:}
\item Si l'évènement $A_{n}$ est réalisé alors l'évènement $A_{n+1}$ est réalisé avec une probabilité 0,3.
\item Si l'évènement $A_{n}$ n'est pas réalisé alors l'évènement $A_{n+1}$ est réalisé avec une probabilité 0,7.
\end{list}

On suppose que $p_{1}=1$.

\subsection*{Partie A :}

\begin{enumerate}
\item On complète l'arbre des probabilités ci-dessous:

\begin{center}
\bigskip
\pstree[treemode=R,nodesep=4pt,levelsep=2.5cm,treesep=1cm]{\TR{$A_1$}}
{
	\pstree{\TR{$A_2$}\naput{\blue $0,3$}}
	{
		\TR{$A_3$}\naput{\blue $0,3$}
		\TR{$\overline{A_3}$}\nbput{\blue $0,7$}
	}
	\pstree{\TR{$\overline{A_2}$}\nbput{$0,7$}}
	{
		\TR{$A_3$}\naput{\blue $0,7$}
		\TR{$\overline{A_3}$}\nbput{\blue $0,3$}
	}
}
\bigskip
\end{center}

\item 
$\left \{ A_2\;,\; \overline{A_2} \right \}$ forme une partition donc, d'après la formule des probabilités totales:

$\begin{aligned}
P\left (A_3\right )
& = P\left (A_2 \cap A_3\right ) +P\left (\overline{A_2} \cap A_3\right )
= P\left (A_2\right ) \times P_{A_2}\left (A_3\right ) + P\left (\overline{ A_2} \right ) \times P_{\overline{ A_2} }\left (A_3\right )\\
& = 0,3\times 0,3+0,7\times 0,7 = 0,58
\end{aligned}$

\item 
$P_{A_{3}}\left(A_{2}\right) = \dfrac{P\left (A_2 \cap A_3\right )}{P\left ( A_3\right )}
= \dfrac{0,3\times 0,3}{0,58}\approx 0,16$
\end{enumerate}

\subsection*{Partie B :}

	Dans cette partie, on étudie la suite $\big(p_{n}\big)$ avec $n \geqslant 1$.

\begin{enumerate}
\item On complète l'arbre des probabilités ci-dessous :

\begin{center}
\bigskip
\pstree[treemode=R,nodesepA=0pt,nodesepB=4pt,levelsep=2.5cm,treesep=1cm]{\TR{}}
{
\pstree[nodesepA=4pt]{\TR{$A_{n}$}\naput{$p_n$}}
	{
	\TR{$A_{n+1}$}\naput{\blue $0,3$}
	\TR{$\overline{A_{n+1}}$}\nbput{\blue $0,7$}
	}
\pstree[nodesepA=4pt]{\TR{$\overline{A_{n}}$}\nbput{\blue $1-p_n$}}
	{
	\TR{$A_{n+1}$}\naput{\blue $0,7$}
	\TR{$\overline{A_{n+1}}$}\nbput{\blue $0,3$}
	}
}
\bigskip
\end{center}
	
	\item \begin{enumerate}
		\item 
$\left \{ A_{n}\;,\; \overline{A_n} \right \}$ forme une partition donc, d'après la formule des probabilités totales:

$\begin{aligned}
p_{n+1} = P\left (A_{n+1}\right )
&  = P\left (A_n \cap A_{n+1}\right ) +P\left (\overline{A_n} \cap A_{n+1}\right )\\
& = P\left (A_n\right ) \times P_{A_n}\left (A_{n+1}\right ) + P\left (\overline{ A_n} \right ) \times P_{\overline{ A_n} }\left (A_{n+1}\right )\\
& = p_n\times 0,3+\left (1-p_n\right ) \times 0,7
= 0,3p_n +0,7-0,7p_n = -0,4p_n + 0,7
\end{aligned}$
		\end{enumerate}
	On considère la suite $\left (u_{n}\right )$, définie pour tout entier naturel $n$ non nul par : 
$u_{n}=p_{n}-0,5$. Donc $p_n=u_n+0,5$.

	\begin{enumerate}[resume]
	\item 
$\begin{aligned}[t]
u_{n+1}
&=p_{n+1}-0,5 = -0,4p_n + 0,7-0,5 = -0,4\left (u_n+0,5\right ) +0,2 = -0,4u_n -0,2+0,2\\
&=-0,4u_n
\end{aligned}$

$u_1=p_1-0,5=1-0,5=0,5$

Donc la suite $\left (u_n\right )$ est géométrique de raison $q=-0,4$ et de premier terme $u_1=0,5$.

	\item 
On en déduit que, pour tout $n$, on a:
$u_n=u_1\times q^{n-1}= 0,5 \times \left (-0,4\right )^{n-1}$.

Et donc $p_n=u_n+0,5= 0,5\times \left (-0,4\right )^{n-1} + 0,5$.

\item
On a: $-1 < -0,4 < 1$ donc $\ds\lim_{n\to +\infty} (-0,4)^{n-1} = 0$ donc $\ds\lim_{n\to +\infty} p_n = 0,5$.
	\end{enumerate}
\end{enumerate}

\subsection*{Partie C :}

Soit $x \in ]0 ~;~ 1[$, on suppose que $P_{\overline{A_{n}}}\left(A_{n+1}\right)=P_{A_{n}}\left(\overline{A_{n+1}}\right)=x$. On rappelle que $p_{1}=1$.

On représente cette situation par un arbre de probabilités.

\begin{center}
\bigskip
\pstree[treemode=R,nodesepA=0pt,nodesepB=4pt,levelsep=2.5cm,treesep=1cm]{\TR{}}
{
	\pstree[nodesepA=4pt]{\TR{$A_{n}$}\naput{$p_n$}}
	{
	\TR{$A_{n+1}$}\naput{\blue $1-x$}
	\TR{$\overline{A_{n+1}}$}\nbput{\blue $x$}
	}
	\pstree[nodesepA=4pt]{\TR{$\overline{A_{n}}$}\nbput{\blue $1-p_n$}}
	{
	\TR{$A_{n+1}$}\naput{\blue $x$}
	\TR{$\overline{A_{n+1}}$}\nbput{\blue $1-x$}
	}
}
\bigskip
\end{center}

\begin{enumerate}
\item 
$\left \{ A_{n}\;,\; \overline{A_n} \right \}$ forme une partition donc, d'après la formule des probabilités totales:

$\begin{aligned}
p_{n+1} = P\left (A_{n+1}\right)
&  = P\left (A_n \cap A_{n+1}\right ) +P\left (\overline{A_n} \cap A_{n+1}\right )\\
& = P\left (A_n\right ) \times P_{A_n}\left (A_{n+1}\right ) + P\left (\overline{ A_n} \right ) \times P_{\overline{ A_n} }\left (A_{n+1}\right )\\
& = p_n\times \left (1-x\right )+\left (1-p_n\right ) \times x
= p_n - x p_n + x - xp_n = \left (1-2x\right ) p_n + x
\end{aligned}$

\item On va démontrer par récurrence que $p_{n}=\frac{1}{2}(1-2 x)^{n-1}+\frac{1}{2}$ pour tout $n\geqslant 1$.

\begin{list}{\textbullet}{}
\item \textbf{Initialisation}

Pour $n=1$, on a $\frac{1}{2}(1-2 x)^{1-1}+\frac{1}{2}=\frac{1}{2}+\frac{1}{2}=1$

Or $p_1=1$, donc la propriété est vraie pour $n=1$.
\item \textbf{Hérédité}

On suppose la propriété vraie au rang $n$ c'est-à-dire que $p_{n}=\frac{1}{2}(1-2 x)^{n-1}+\frac{1}{2}$; c'est l'hypothèse de récurrence.

$\begin{aligned}
p_{n+1}
& = \left (1-2x\right )p_n+x 
= \left (1-2x\right ) \left ( \frac{1}{2}(1 - 2 x)^{n-1}+\frac{1}{2}\right )+x\\
& = \left (1-2x\right ) \times \frac{1}{2}(1 - 2x)^{n-1}+ \left (1 - 2x\right ) \times\frac{1}{2}+x
 = \frac{1}{2}(1 - 2x)^{n} +\frac{1}{2} - x + x\\
& = \frac{1}{2}(1 - 2x)^{n} +\frac{1}{2}
\end{aligned}$

La propriété est donc vraie au rang $n+1$.
\item \textbf{Conclusion}

La propriété est vraie au rang 1 et elle est héréditaire pour tout $n\geqslant 1$; d'après le principe de récurrence, la propriété est vraie pour tout entier $n$ non nul.
\end{list}

On a donc démontré par récurrence sur $n$ que $p_{n}=\frac{1}{2}(1-2 x)^{n-1}+\frac{1}{2}$ pour tout $n\geqslant 1$.

\item 
$0 < x < 1$ donc $-2 < -2x < 0$ et donc $-1 < 1 - 2x < 1$.

On en déduit que $\lim_{n\to +\infty} (1-2x)^{n-1}=0$ 
donc que $\lim_{n\to +\infty} p_n=\frac{1}{2}$.
\end{enumerate}

\end{document}
